% preamble.tex — shared by every chapter of hol-kanren.tex; tufte-book class, compile with XeLaTeX
% (tectonic hol-kanren.tex, or latexmk -xelatex). The class is loaded in hol-kanren.tex as
%   \PassOptionsToPackage{dvipsnames}{xcolor}
%   \documentclass[a4paper,justified,nofonts,nobib,nols]{tufte-book}
% tufte-book already loads fontspec (through xltxtra), xcolor, hyperref, titlesec, changepage and geometry.
%
% Layout: a 107mm text column and a 49.4mm margin column. The margin holds \sidenote, \marginnote and the
% captions of figure/table floats. Code listings and the reading boxes span the full width (text + margin),
% so that code lines keep their length; wide tables and rule figures use fullwidth, figure* or table*.
\usepackage{amsmath,amsthm}
\usepackage{unicode-math}
\setmainfont{LibertinusSerif}[Extension=.otf,UprightFont=*-Regular,BoldFont=*-Bold,ItalicFont=*-Italic,BoldItalicFont=*-BoldItalic]
\setsansfont{LibertinusSans}[Extension=.otf,UprightFont=*-Regular,BoldFont=*-Bold,ItalicFont=*-Italic]
\setmathfont{LibertinusMath-Regular.otf}
\setmonofont{DejaVu Sans Mono}[Scale=MatchLowercase]
\usepackage{mathpartir}
\usepackage[most]{tcolorbox}
\usepackage{booktabs,longtable,array,tabularx}
\usepackage{enumitem}
\usepackage{csquotes}
\usepackage{microtype}
\hypersetup{hidelinks}
\usepackage[capitalise,nameinlink]{cleveref}

% ---------------------------------------------------------------- tufte adjustments
% Number chapters, sections and subsections (tufte numbers nothing), so that \cref targets read "Section 9.8".
\setcounter{secnumdepth}{2}
\setcounter{tocdepth}{1}
% The width of text + margin, and the part of it that lies in the margin.
\newlength{\fullwidthlen}
\setlength{\fullwidthlen}{\textwidth}
\addtolength{\fullwidthlen}{\marginparsep}
\addtolength{\fullwidthlen}{\marginparwidth}
\newlength{\overhang}
\setlength{\overhang}{\marginparsep}
\addtolength{\overhang}{\marginparwidth}
% A longtable starts at the left edge and may run into the margin: give its columns widths in \fullwidthlen.
\setlength{\LTleft}{0pt}
\setlength{\LTright}{0pt plus 1fill minus \overhang}
% Inside a reading box, code and tables keep to the box (which is already full width).
\newif\ifinbox \inboxfalse

% ---------------------------------------------------------------- code
% fvextra, not listings: listings scrambles code points above 255 (λ, α, ∀, ...) under XeTeX.
% Environments: lprolog (λProlog source), scheme (aux Scheme source), answers (actual output of a run);
% each takes optional fvextra settings, e.g. \begin{scheme}[fontsize=\footnotesize]. Outside a box they
% span the full width. Inline: \lp{...} and \scm{...} (balanced braces; not inside section titles or
% captions, use \texttt there).
\usepackage{fvextra}
\fvset{fontsize=\small,breaklines=true,frame=leftline,framerule=1.2pt,xleftmargin=1.2em,framesep=0.6em,tabsize=2}
\newcommand{\codewidebegin}{\ifinbox\else\fvset{xrightmargin=-\overhang}\fi}% local to the environment
\newcommand{\codewideend}{}
\DefineVerbatimEnvironment{lprologV}{Verbatim}{rulecolor=\color{Plum!60}}
\DefineVerbatimEnvironment{schemeV}{Verbatim}{rulecolor=\color{OliveGreen!60}}
\DefineVerbatimEnvironment{answersV}{Verbatim}{fontsize=\footnotesize,formatcom=\color{black!75},rulecolor=\color{Gray!30}}
% \VerbatimEnvironment comes first, so that fancyvrb waits for \end{scheme}, not \end{fullwidth}; the inner environment
% parses the optional [...] itself (a wrapper with its own optional argument would eat the line end).
\newenvironment{lprolog}{\VerbatimEnvironment\codewidebegin\begin{lprologV}}{\end{lprologV}\codewideend}
\newenvironment{scheme}{\VerbatimEnvironment\codewidebegin\begin{schemeV}}{\end{schemeV}\codewideend}
\newenvironment{answers}{\VerbatimEnvironment\codewidebegin\begin{answersV}}{\end{answersV}\codewideend}
\newcommand{\lp}{\Verb[breaklines=true,fontsize=\small,formatcom=\color{Plum!70!black}]}
\newcommand{\scm}{\Verb[breaklines=true,fontsize=\small,formatcom=\color{OliveGreen!70!black}]}

% ---------------------------------------------------------------- theorem-like environments
\theoremstyle{definition}
\newtheorem{definition}{Definition}[chapter]
\newtheorem{example}[definition]{Example}
\theoremstyle{plain}
\newtheorem{theorem}[definition]{Theorem}
\newtheorem{proposition}[definition]{Proposition}
\newtheorem{lemma}[definition]{Lemma}
\newtheorem{corollary}[definition]{Corollary}
\theoremstyle{remark}
\newtheorem{remark}[definition]{Remark}

% ---------------------------------------------------------------- reading boxes
% tutorial: informal walk-through; kanren: how the concept maps to (aux kanren micro); pitfall: a trap the
% book stresses. They span the full width; inside them code and tables keep to the box. No \sidenote or
% \marginnote inside a box (the margin is taken by the box).
\tcbset{hkbox/.style={enhanced,breakable,grow to right by=\overhang,fonttitle=\bfseries\sffamily,
  boxrule=0.6pt,arc=1.5pt,before upper={\inboxtrue}}}
\newtcolorbox{tutorial}[1][]{hkbox,colback=blue!3,colframe=blue!45!black,
  title={Tutorial\ifx\relax#1\relax\else: #1\fi}}
\newtcolorbox{kanren}[1][]{hkbox,colback=green!3,colframe=green!35!black,
  title={Kanren angle\ifx\relax#1\relax\else: #1\fi}}
\newtcolorbox{pitfall}[1][]{hkbox,colback=orange!4,colframe=orange!60!black,
  title={Pitfall\ifx\relax#1\relax\else: #1\fi}}

% ---------------------------------------------------------------- notation
\newcommand{\book}[1]{\textup{[MN~p.\,#1]}}          % a PRINTED page of Miller & Nadathur (= PDF page - 16)
\newcommand{\imp}{\supset}                              % implication, as in the book
\newcommand{\seq}[3]{#1;#2\longrightarrow#3}            % Σ; P ⟶ G, a sequent of uniform proof search
\newcommand{\Sig}{\Sigma}
\newcommand{\Prog}{\mathcal{P}}
\newcommand{\Llam}{L_\lambda}
\newcommand{\fohc}{\textit{fohc}}  \newcommand{\fohh}{\textit{fohh}}
\newcommand{\hohc}{\textit{hohc}}  \newcommand{\hohh}{\textit{hohh}}
\newcommand{\typ}[1]{\mathsf{#1}}                        % a type or sort, e.g. \typ{o}, \typ{int}
\newcommand{\cst}[1]{\mathsf{#1}}                        % a constant of a signature
\newcommand{\abs}[2]{\lambda #1.\,#2}
\newcommand{\eqb}{=_{\beta\eta}}
\newcommand{\subst}[3]{#1[#2/#3]}                        % t[s/x]
\newcommand{\rulename}[1]{\textsc{#1}}
\newcommand{\aux}{\textsf{aux}}
\newcommand{\micro}{\textsf{(aux kanren micro)}}
